\section*{Capítulo \ref{ch:b1:stereochemistry-in-depth} --- Estereoquímica en profundidad}

\begin{solution}{exo:b1:stereochemistry-in-depth:1}
\omterm{def:b1:stereochemistry-in-depth:isomers}{Isómeros constitucionales}; \omterm{def:b1:stereochemistry-in-depth:enantiomers}{enantiómeros}; \omterm{def:b1:stereochemistry-in-depth:enantiomers}{diastereoisómeros}
(\omterm{def:b1:stereochemistry-in-depth:isomers}{estereoisómeros} que no son imágenes especulares); dos \omterm{def:b1:stereochemistry-in-depth:isomers}{conformaciones} de
un mismo compuesto; \omterm{def:b1:stereochemistry-in-depth:enantiomers}{diastereoisómeros} (\cip{R,S} es la forma meso).
\end{solution}

\begin{solution}{exo:b1:stereochemistry-in-depth:2}
\ce{-OH} $>$ \ce{-NH2} $>$ \ce{-CH3} $>$ \ce{-H} (O, N, C, H).
\ce{-COOH} (O,O,O) $>$ \ce{-CHO} (O,O,H) $>$ \ce{-CH2OH} (O,H,H) $>$
\ce{-CH(CH3)2} (C,C,H). \ce{-Br} $>$ \ce{-Cl} $>$ \ce{-CCl3} (Cl,Cl,Cl)
$>$ \ce{-CH2Cl} (Cl,H,H): el primer átomo decide antes que sus vecinos.
\end{solution}

\begin{solution}{exo:b1:stereochemistry-in-depth:3}
2,3-Diclorobutano: dos centros equivalentes, \cip{R,R}, \cip{S,S} y el
meso \cip{R,S}: tres. 2-Bromo-3-clorobutano: dos centros distintos, cuatro
\omterm{def:b1:stereochemistry-in-depth:isomers}{estereoisómeros}, sin forma meso. Pentano-2,4-diol: tres, siendo meso el
\cip{2R,4S}.
\end{solution}

\begin{solution}{exo:b1:stereochemistry-in-depth:4}
En C5: isopropenilo (C,C,C, pues el enlace doble cuenta dos veces) $>$ C6
$>$ C4 $>$ H. C6 y C4 llevan ambos (C,H,H); un paso más allá, C6 conduce
al carbono carbonílico (O,O,C) y C4 a C3 (C,C,H): gana C6. Con el
isopropenilo delante y H detrás, el orden isopropenilo $\to$ C6 $\to$ C4
gira en el sentido de las agujas del reloj tal como está dibujado:
\cip{R}, la carvona de la hierbabuena; la de la cuña rayada es \cip{S}.
\end{solution}

\begin{solution}{exo:b1:stereochemistry-in-depth:5}
$\theta = 0$: metilos eclipsados, \qty{16.5}{kJ/mol}; $60^\circ$: gauche,
unos \qty{2.9}{kJ/mol} (mínimo de 2.8 en $62^\circ$); $120^\circ$: metilo
eclipsando H, \qty{15.2}{kJ/mol}; $180^\circ$: anti, 0. Con energías
iguales, dos \omterm{def:b1:stereochemistry-in-depth:conformations}{conformaciones gauche} por cada anti: dos tercios gauche.
\end{solution}

\begin{solution}{exo:b1:stereochemistry-in-depth:6}
Cadena vertical con \ce{COOH} arriba y \ce{CH3} abajo; prioridades
$\ce{NH2} > \ce{COOH} > \ce{CH3} > \ce{H}$. Poniendo \ce{NH2} a la
izquierda y H a la derecha, $\ce{NH2} \to \ce{COOH} \to \ce{CH3}$ gira en
el sentido de las agujas del reloj tal como se ve; H es horizontal (hacia
el lector), de modo que el descriptor se invierte: \cip{S}. El grupo
amino está a la izquierda.
\end{solution}

\begin{solution}{exo:b1:stereochemistry-in-depth:7}
\textit{cis}: en carbonos contiguos, un enlace hacia arriba y otro hacia
abajo entre las posiciones \omterm{def:b1:stereochemistry-in-depth:conformations}{axiales} y \omterm{def:b1:stereochemistry-in-depth:conformations}{ecuatoriales} da axial--ecuatorial en
las dos \omterm{def:b1:stereochemistry-in-depth:conformations}{sillas}. \textit{trans}: diecuatorial o diaxial. En la \omterm{def:b1:stereochemistry-in-depth:conformations}{silla}
diecuatorial, los dos metilos están en gauche entre sí (una interacción);
el isómero \textit{cis} tiene esa interacción más un metilo \omterm{def:b1:stereochemistry-in-depth:conformations}{axial}, unos
\qty{5.7}{kJ/mol} más. El isómero \textit{trans} es el más estable.
\end{solution}

\begin{solution}{exo:b1:stereochemistry-in-depth:8}
$[\alpha] = 13.2/(2.00 \times 0.100) = +66.0^\circ$. Un tubo de
\qty{1.00}{dm}: $+6.60^\circ$. El enantiómero: $-13.2^\circ$ en el tubo de
\qty{2.00}{dm}.
\end{solution}

\begin{solution}{exo:b1:stereochemistry-in-depth:9}
\qty{25}{\celsius}: $\exp(5700/(8.314 \times 298.15)) = 10.0$, fracción
$10.0/11.0 = \qty{91}{\percent}$. \qty{-50}{\celsius}: $\exp(5700/(8.314
\times 223.15)) = 21.6$, \qty{96}{\percent}.
\end{solution}

\begin{solution}{exo:b1:stereochemistry-in-depth:10}
En el etano, los seis hidrógenos son iguales: todas las \omterm{def:b1:stereochemistry-in-depth:conformations}{conformaciones
alternadas} son la misma. En el butano, el metilo trasero puede situarse
opuesto al delantero (anti) o a su lado ($\pm 60^\circ$, gauche), donde
los dos metilos se estorban. Poblaciones: anti 1, gauche
$2\exp(-2.8/2.48) = 0.65$: alrededor de un \qty{61}{\percent} anti y un
\qty{39}{\percent} gauche.
\end{solution}

\begin{solution}{exo:b1:stereochemistry-in-depth:11}
C2 y C4 son estereogénicos. C3 lleva H, OH y dos ramas idénticas en
constitución; solo difieren cuando C2 y C4 tienen descriptores opuestos,
y solo entonces es C3 estereogénico (pseudoasimétrico). Isómeros:
\cip{2R,4R} y \cip{2S,4S}, una pareja de \omterm{def:b1:stereochemistry-in-depth:enantiomers}{enantiómeros} (C3 no
estereogénico); \cip{2R,4S} con las dos disposiciones en C3, dos formas
meso. Cuatro en total.
\end{solution}

\begin{solution}{exo:b1:stereochemistry-in-depth:12}
$[\alpha]_{\text{muestra}} = 2.31/(1.00 \times 0.050) = +46.2^\circ$, y
$2x - 1 = 46.2/66.0 = 0.700$: $x = 0.850$, un \qty{85}{\percent} del
enantiómero $(+)$ y un \qty{15}{\percent} del $(-)$.
\end{solution}

\begin{solution}{pb:b1:stereochemistry-in-depth:1}
\textbf{1.} Anillo de ciclohexano: C1 lleva el OH; C2, el grupo
isopropilo, y C5, el metilo; C1, C2 y C5 son estereogénicos.
\textbf{2.} $2^3 = 8$, cuatro parejas de \omterm{def:b1:stereochemistry-in-depth:enantiomers}{enantiómeros}.
\textbf{3.} Ninguna disposición tiene un plano especular: los tres
sustituyentes son todos distintos.
\textbf{4.} \ce{OH} $>$ C2 (C,C,H) $>$ C6 (C,H,H) $>$ H.
\textbf{5.} C1 (O,C,H) $>$ isopropilo (C,C,H) $>$ C3 (C,H,H) $>$ H.
\textbf{6.} C6 y C4, ambos (C,H,H); a continuación, C6 conduce a C1
(O,C,H), y C4, a C3
(C,H,H): C6 $>$ C4 $>$ \ce{CH3} $>$ H.
\textbf{7.} \cip{1S,2R,5S}.
\textbf{8.} En una \omterm{def:b1:stereochemistry-in-depth:conformations}{silla}, dos enlaces \omterm{def:b1:stereochemistry-in-depth:conformations}{ecuatoriales} en carbonos contiguos
(C1, C2) apuntan uno hacia arriba y otro hacia abajo: \textit{trans}; en
los carbonos 1 y 3 del anillo (C1 y C5, a través de C6), los dos enlaces
\omterm{def:b1:stereochemistry-in-depth:conformations}{ecuatoriales} apuntan en el mismo sentido: \textit{cis}. Esa es la
disposición del mentol.
\textbf{9.} Los tres pasan a ser \omterm{def:b1:stereochemistry-in-depth:conformations}{axiales}.
\textbf{10.} En la \omterm{def:b1:stereochemistry-in-depth:conformations}{silla} invertida, solo el metilo costaría unos
\qty{5.7}{kJ/mol}, y el grupo isopropilo, mayor, más; el OH añade un coste
propio menor, y el OH y el metilo \omterm{def:b1:stereochemistry-in-depth:conformations}{axiales}, en los carbonos 1 y 3 y en la
misma cara, se estorbarían además entre sí. Eso supera con creces
\qty{12}{kJ/mol}, una razón de más de $\exp(12000/2479) \approx 130$ a
uno: el mentol está casi por completo en la \omterm{def:b1:stereochemistry-in-depth:conformations}{silla} totalmente \omterm{def:b1:stereochemistry-in-depth:conformations}{ecuatorial}.
\textbf{11.} Un diastereoisómero: solo difiere uno de los tres centros.
\textbf{12.} Con el isopropilo y el metilo \omterm{def:b1:stereochemistry-in-depth:conformations}{ecuatoriales}, el OH es \omterm{def:b1:stereochemistry-in-depth:conformations}{axial}.
\textbf{13.} Los \omterm{def:b1:stereochemistry-in-depth:enantiomers}{diastereoisómeros} difieren en todas sus propiedades (aquí,
en los \omterm{def:b1:intermolecular-forces-solvents:hbond}{enlaces de hidrógeno} de un OH \omterm{def:b1:stereochemistry-in-depth:conformations}{axial} o \omterm{def:b1:stereochemistry-in-depth:conformations}{ecuatorial}); los
\omterm{def:b1:stereochemistry-in-depth:enantiomers}{enantiómeros} solo difieren frente a participantes \omterm{def:b1:stereochemistry-in-depth:stereocentre}{quirales}.
\textbf{14.} $[\alpha] = -2.46/(1.00 \times 0.0500) = -49.2^\circ$, dentro
del intervalo del manual.
% ledger: alpha:l-menthol-range
\textbf{15.} $+49.2^\circ$; $0$.
\textbf{16.} $-1.97/0.0500 = -39.4^\circ$.
\textbf{17.} $\alpha = \ell c[\alpha]_{(-)}\bigl(x - (1 - x)\bigr) =
\ell c[\alpha]_{(-)}(2x - 1)$.
\textbf{18.} $x = \frac12(1 + \alpha/\alpha_{\text{ref}})$ con $\ell$ y $c$
iguales.
\textbf{19.} Sí: cualquier otro compuesto ópticamente activo añade su
propio término (\omterm{prop:b1:stereochemistry-in-depth:biot}{ley de Biot}). Hace falta comprobar la pureza química por
otro método (cromatografía) antes de interpretar la rotación como una
composición.
\textbf{20.} $x = \frac12(1 + 1.97/2.46) = 0.900$.
\textbf{21.} $x = \frac12(1 + 2.40/2.46) = 0.988$.
\textbf{22.} La segunda (\qty{98.8}{\percent}).
\textbf{23.} $-3.94^\circ$; el mismo error de lectura es entonces una
fracción menor del ángulo.
\textbf{24.} \textbf{\qty{90.0}{\percent} de $(-)$-mentol} (y un
\qty{10.0}{\percent} de su enantiómero).
\end{solution}
